\section*{Chapitre \ref{ch:m2:swap-clearing-and-the-ccp-basis} --- Compensation des swaps et base entre chambres de compensation}

\begin{solution}{exo:m2:swap-clearing-and-the-ccp-basis:1}
$36.4 \times 81\,109 = \text{USD}~2.95$ millions~: 36.4 points de base de taux, appliqués
au DV01.
\end{solution}

\begin{solution}{exo:m2:swap-clearing-and-the-ccp-basis:2}
Un receveur perd lorsque les taux montent~: $3 \times 81\,109 = \text{USD}~243\,327$, versé en espèces
par l'intermédiaire de son courtier compensateur à la chambre le jour même ou le lendemain matin.
\end{solution}

\begin{solution}{exo:m2:swap-clearing-and-the-ccp-basis:3}
Si le \omterm{def:m2:swap-clearing-and-the-ccp-basis:mandate}{portage des positions} réussit, les swaps du client et sa marge passent à un autre
membre compensateur et se poursuivent sans modification~: ses couvertures subsistent. En cas d'échec, la
chambre clôture les positions, en restituant la marge du client nette des pertes~;
le client doit reconstituer ses couvertures sur le marché alors en vigueur, après
avoir été sans couverture dans l'intervalle.
\end{solution}

\begin{solution}{exo:m2:swap-clearing-and-the-ccp-basis:4}
USD~76\,442 sur la durée de vie du swap~; divisé par le DV01 de 81\,109, une base
courante de 0.94 points de base par an.
\end{solution}

\begin{solution}{exo:m2:swap-clearing-and-the-ccp-basis:5}
À la chambre où les clients paient le taux fixe. Les dealers qui y reçoivent se retrouvent
à sens unique et doivent se couvrir ailleurs, en déposant deux fois la marge~; ils demandent un taux fixe plus élevé pour
recevoir, si bien que le \omterm{def:m2:government-bonds:zero}{taux au pair} y est plus élevé. L'analyse de 2017 a relevé des swaps
cotés au CME au-dessus des mêmes swaps au LCH, pour exactement cette raison.
\end{solution}

\begin{solution}{exo:m2:swap-clearing-and-the-ccp-basis:6}
Oui~: 9~milliards d'euros dépassent le seuil de 8~milliards d'euros de la phase finale.
La marge initiale n'est échangée qu'au-delà du seuil de 50~millions d'euros par
groupe de contreparties~: 20~millions d'euros, si rien d'autre dans le groupe n'utilise le
seuil.
\end{solution}

\begin{solution}{exo:m2:swap-clearing-and-the-ccp-basis:7}
1.15, 2.30 et 3.45 points de base par position à sens unique. À 50, deux chambres
requièrent 4.60, contre les 3.4 de l'exemple de 2017~: le même ordre de grandeur, la
différence tenant à la volatilité illustrative, à la période de marge et au
spread de financement, et aux modèles de marge réels qui compensent les positions du portefeuille d'un dealer.
\end{solution}

\begin{solution}{exo:m2:swap-clearing-and-the-ccp-basis:8}
Les deux swaps se compensent en risque mais pas en marge~: chaque chambre exige la marge
de son côté, si bien que l'«~arbitrage~» dépose deux marges initiales pendant trente ans. Au
coût de financement du marché, c'est à cela que servent les 3 points de base. Seule une firme
dont les positions existantes font que chaque nouvelle transaction \emph{réduit} sa marge aux deux
chambres gagne le spread, et c'est le créneau que se disputent les dealers.
\end{solution}

\begin{solution}{pb:m2:swap-clearing-and-the-ccp-basis:1}
\textbf{1.} $2 \times 10^9 \times 8.111 \times 10^{-4} = \text{USD}~1.62$ millions par point de base.
\textbf{2.} $36.4 \times 1.622 \approx \text{USD}~59.1$ millions.
\textbf{3.} $5 \times 1.622 = \text{USD}~8.11$ millions, versés par le dealer qui reçoit
le taux fixe en A~; en B, où il paie le taux fixe, il reçoit le même montant.
\textbf{4.} Chaque chambre ne garantit que ses propres transactions et doit pouvoir clôturer
le portefeuille d'un défaillant \emph{sur place}~; une couverture chez une autre chambre ne
réduit pas sa perte.
\textbf{5.} Elle tendrait vers zéro~: une position à plat en A n'a aucun risque de marché
là-bas, et seulement une petite marge au titre de la base et de la liquidité.
\textbf{6.} Les mêmes USD~59.1 millions.
\textbf{7.} USD~1.53 millions chacune, USD~3.06 millions pour les deux.
\textbf{8.} 0.94 points de base pour l'une, 1.88 pour les deux.
\textbf{9.} 4.60 points de base pour les deux.
\textbf{10.} $\sqrt 2$ fois autant~: USD~83.5 millions dans chaque chambre.
\textbf{11.} En payant le taux fixe en A à un dealer qui y reçoit, sa position en
A est à plat~: pas de marge initiale notable en A, ni en B, où elle n'a plus
besoin de couverture.
\textbf{12.} Jusqu'à 1.88 points de base de taux~: ce que coûteraient les deux marges
à sens unique sur les dix ans.
\textbf{13.} Parce que chaque dealer en A a affaire aux mêmes clients~: tous sont receveurs
là-bas, et aucun ne veut aggraver le déséquilibre en recevant davantage.
\textbf{14.} Dans le taux fixe qu'ils paient en A, supérieur à celui de B d'environ le coût
des deux marges des dealers.
\textbf{15.} Des clients qui reçoivent le taux fixe en A (fonds de pension, assureurs), une règle
qui déplace des flux entre chambres, ou une modification du modèle de marge de l'une des deux
chambres, ou une compensation croisée des marges avec l'autre.
\textbf{16.} Parce qu'il doit être financé~: la chambre rémunère le collatéral en espèces moins
que ce que la firme paie pour l'emprunter, et les titres déposés doivent être financés en
repo~; l'écart est le spread de financement.
\textbf{17.} Elles recourent à la simulation historique de portefeuilles entiers, à des majorations pour
la liquidité et la concentration, et à des planchers~; la compensation au sein d'un portefeuille fait
que la marge incrémentale d'une transaction dépend de tout le reste que détient le membre.
\textbf{18.} Elle force des flux vers une chambre sans égard pour son portefeuille, créant
des books à sens unique là où il n'y en avait pas, et faisant ainsi naître ou élargir une base
entre cette chambre et les autres.
\textbf{19.} Résultat nommé~: \emph{la base de seuil de rentabilité} pour les swaps à dix ans est de 1.88
points de base~: le coût courant de la marge des positions à sens unique dans les deux chambres, que
le dealer devrait accepter de payer pour compenser en A à la place.
\textbf{20.} Parce que la marge est calculée séparément dans chaque chambre de compensation, et
le coût de détention d'un swap dépend de ceux qui y négocient également.
\end{solution}

\begin{solution}{iq:m2:swap-clearing-and-the-ccp-basis:1}
Les swaps bilatéraux avaient créé un réseau d'expositions entre les grandes banques, opaque aux
régulateurs et aux autres, ce qui amplifiait en 2008 la crainte de la défaillance d'un seul.
Le G20 a exigé que les swaps standardisés soient compensés, déclarés et, lorsque
c'est possible, négociés sur des plateformes. Pour un dealer~: sa contrepartie sur la plupart des swaps
est une chambre de compensation~; il dépose chaque jour une marge initiale et une marge de variation~; son
bilan profite de la compensation~; ses coûts incluent désormais le financement des marges et
les contributions au fonds de défaillance.
\iqlookfor{le réseau de contreparties comme problème et la marge comme nouveau coût.}
\end{solution}

\begin{solution}{iq:m2:swap-clearing-and-the-ccp-basis:2}
La marge de variation règle la variation quotidienne de la valeur d'une position, de sorte que
aucune partie n'accumule de gain ou de perte non réglé. La marge initiale est un collatéral
déposé d'avance et conservé, pour couvrir la perte que la chambre ou la contrepartie pourrait
subir en clôturant les positions d'un défaillant pendant la période de marge en
risque~; elle est restituée lorsque la position est clôturée.
\iqlookfor{règlement-livraison des variations passées, opposé au coussin pour les variations futures.}
\end{solution}

\begin{solution}{iq:m2:swap-clearing-and-the-ccp-basis:3}
L'écart de \omterm{def:m2:government-bonds:zero}{taux au pair} entre des swaps identiques compensés chez CME et chez LCH.
Les flux de clients sont unidirectionnels dans une chambre donnée, de sorte que les dealers se retrouvent du même côté
dans chaque chambre et déposent deux marges~; la base tarife ce coût de financement. Il ne s'agit pas
d'un arbitrage~: la capturer recrée les deux marges qu'elle rémunère, sauf si les positions existantes
font des transactions des opérations réduisant la marge.
\iqlookfor{le déséquilibre des flux et le mécanisme de marge, et pourquoi ce n'est pas
de l'argent gratuit.}
\end{solution}

\begin{solution}{iq:m2:swap-clearing-and-the-ccp-basis:4}
L'ajustement de valorisation de marge (MVA)~: la valeur présente du financement de la marge initiale exigée par
une transaction sur toute sa durée de vie. L'estimer en projetant dans le temps la marge initiale incrémentale
de la transaction (à partir de son profil de DV01 futur, ou en simulant la marge du portefeuille), en multipliant
par le spread de financement puis en actualisant~; l'exprimer sous forme de spread courant en divisant par le DV01
de la transaction.
\iqlookfor{marge incrémentale et non autonome, et une projection dans
le temps.}
\end{solution}

\begin{solution}{iq:m2:swap-clearing-and-the-ccp-basis:5}
Bilatéral~: le dealer est la contrepartie, si bien que le risque de crédit est géré par les clauses de la
\omterm{def:m2:swap-clearing-and-the-ccp-basis:umr}{credit support annex} (CSA)~; si les deux parties dépassent les seuils, la marge initiale est échangée
et ségrégée, en général plus élevée qu'en chambre de compensation car la période de marge est plus longue
(dix jours au lieu de cinq) et le netting plus étroit~; le dealer facture le risque de crédit et de
financement (les ajustements de valorisation traités dans le livre~6). Le client gagne
en souplesse sur les conditions et évite les frais de compensation et l'exposition au fonds de défaillance
par l'intermédiaire de son courtier.
\iqlookfor{différences de crédit et de marge, et leur valorisation.}
\end{solution}

\begin{solution}{iq:m2:swap-clearing-and-the-ccp-basis:6}
Répliquer la méthodologie de marge publiée par la chambre~: scénarios historiques
appliqués au portefeuille du membre, avec ses majorations de liquidité et de concentration.
Précalculer chaque matin le vecteur des P\&L en scénarios du portefeuille~; pour une nouvelle
transaction, calculer son vecteur de P\&L en scénarios (linéaire en sensibilités pour aller vite, réévaluation complète
pour les transactions importantes), l'ajouter et lire le nouveau quantile~; l'écart
constitue la marge incrémentale. Conserver les vecteurs en mémoire, les mettre à jour à chaque
transaction, et le confronter au propre calculateur de marge de la chambre.
\iqlookfor{vecteurs de scénarios, quantiles incrémentaux et validation face à la
chambre.}
\end{solution}
